\chapter{Коли шукати, а коли вирішувати: додаємо \nt{NORA}}
\chaptermark{Додаємо \nt{NORA}}
\label{sec:nora}

\section{Чого бракує}

У розділі~\ref{sec:dofa} мережа навчалася завдяки шуму: випадкові відхилення активності були пошуком, а \nt{DOPA} повідомляв, чи стало від них краще. Але там само лишився незаданим один параметр --- \emph{скільки} шукати. Якщо шум надто сильний, мережа метається і не користується тим, що вже знає. Якщо надто слабкий, вона раз у раз вибирає те саме і не помічає, що десь поруч стало краще. У навчанні з підкріпленням це називають вибором між використанням і пошуком, і кожен алгоритм має для нього ручку. У розділі~\ref{sec:neurontypes} ми припустили, що в мозку цю ручку крутить \nt{NORA}.

Норадреналінових нейронів у людини кілька десятків тисяч (таблиця~\ref{tab:types}), майже всі вони зібрані в одній маленькій групі, а їхні виходи розходяться практично по всьому мозку, хоча різні частини цієї групи, як з'ясувалося, частково обслуговують різні області~\cite{Poe2020}. Це широкомовний канал, ще вужчий, ніж \nt{DOPA}. Працює він у двох режимах~\cite{AstonJones2005}. У \emph{фоновому} нейрони видають імпульси рівно, з частотою від часток до кількох імпульсів на секунду, і ця частота повільно змінюється разом із рівнем неспання. У \emph{уривчастому} вони відповідають короткою пачкою на важливу для поточної задачі подію, приблизно в момент ухвалення рішення. Зручна, але неточна контрольна точка --- діаметр зіниці: він змінюється разом з активністю цих нейронів, але також разом з активністю інших областей~\cite{Joshi2016} і холінергічних виходів у корі~\cite{Reimer2016}. Зіниця показує загальний рівень збудження, а не лише \nt{NORA}.

\section{Коефіцієнт підсилення}

Виміряно ось що: норадреналін пригнічує фонову активність нейронів сильніше, ніж їхню відповідь на стимул, і відношення сигналу до фону зростає~\cite{Foote1975NE}. Що крутість характеристики окремого нейрона при цьому буквально множиться на коефіцієнт, з досліду не випливає. Як і в розділі~\ref{sec:neurontypes}, візьмімо просту модель, яка передає цей ефект: норадреналін змінює крутість передавальної характеристики отримувача~\cite{ServanSchreiber1990}. Слабкі, підпорогові входи (фон) тоді дають ще менше, а сильні --- ще більше. Нехай нейрон відповідає частотою
\begin{empheq}[box=\keybox]{equation}
r \;=\; F\bigl(g\,(u-\theta)\bigr), \qquad F(z)=\frac{r_{\max}}{1+e^{-z}} ,
\label{eq:gain}
\end{empheq}
де $u$ --- вхідний струм, $\theta$ --- поріг, а $g$ --- коефіцієнт підсилення, яким у моделі керує \nt{NORA}. За малого $g$ частота плавно йде за входом у широкому діапазоні: нейрон --- \emph{аналоговий підсилювач}. За великого $g$ майже будь-який вхід вище порога дає повну частоту, а нижче --- нуль: нейрон --- \emph{компаратор}, майже цифровий елемент (рис.~\ref{fig:gaincurves}). Одна ручка перемикає той самий нейрон між двома режимами, які в електроніці виконують різні схеми.

\begin{figure}[t]
\centering
\begin{tikzpicture}[>=Stealth,x=1.1cm,y=2.4cm]
\draw[->,gray!70!black] (-3.3,0) -- (3.4,0) node[below left,black] {\small $u$};
\draw[->,gray!70!black] (-3.3,0) -- (-3.3,1.2) node[left,black] {\small $r$};
\draw[densely dashed,gray!60!black] (0,0) -- (0,1.1);
\node[below,font=\small] at (0,0) {$\theta$};
\draw[densely dotted,gray!60!black] (-3.3,1) -- (3.2,1) node[right,black,font=\small] {$r_{\max}$};
\draw[very thick,blue!60!black] plot[domain=-3.2:3.2,samples=60] (\x,{1/(1+exp(-0.8*\x))});
\draw[very thick,orange!85!black] plot[domain=-3.2:3.2,samples=60] (\x,{1/(1+exp(-2*\x))});
\draw[very thick,red!70!black] plot[domain=-3.2:3.2,samples=120] (\x,{1/(1+exp(min(-8*\x,9)))});
\node[blue!60!black,font=\small,anchor=west] at (0.5,0.12) {мале $g$: підсилювач};
\node[red!70!black,font=\small,anchor=east] at (-0.3,0.85) {велике $g$: компаратор};
\end{tikzpicture}
\caption{Передавальна характеристика~\eqref{eq:gain} за трьох значень коефіцієнта підсилення $g$.}
\label{fig:gaincurves}
\end{figure}

Чи допомагає велике підсилення проти шуму? Не завжди. Нехай два входи відрізняються на $\Delta u$, і треба сказати, який більший. Якщо шум $\sigma_{\text{до}}$ доданий до входу \emph{до} підсилювача, то підсилюється і сигнал, і шум, і їхнє відношення не змінюється. Якщо ж шум $\sigma_{\text{після}}$ додається \emph{після} підсилювача --- від ненадійних синапсів і випадкових імпульсів самої мережі, як у розділі~\ref{sec:code}, --- то відношення сигналу до шуму
\begin{empheq}[box=\keybox]{equation}
d' \;=\; \frac{g\,\Delta u}{\sqrt{g^2\sigma_{\text{до}}^2+\sigma_{\text{після}}^2}}
\label{eq:gainsnr}
\end{empheq}
зростає з $g$, доки $g\sigma_{\text{до}}$ не зрівняється з $\sigma_{\text{після}}$~\cite{ServanSchreiber1990}.

\begin{keyblock}
\begin{proposition}[Коли підсилення допомагає]
\label{prop:gainnoise}
Збільшуючи коефіцієнт підсилення, \nt{NORA} покращує розрізнення входів лише проти того шуму, що виникає \emph{після} підсилювача, у самій мережі. Шуму, що прийшов разом із входом, підсилення не прибере.
\end{proposition}
\end{keyblock}

\section{Підсилення у виборі з двох}

Повернімося до вибору з розділу~\ref{sec:gaba}: дві групи \nt{GLUT}-нейронів із входами $I_1$, $I_2$ гальмують одна одну із силою $\beta$. Тепер кожна група має коефіцієнт підсилення $g$: у рівняннях вибору~\eqref{eq:wta} вираз у дужках множиться на $g$. Поки працюють обидві групи, усталені частоти знаходимо з $r_1=g(I_1-\beta r_2)$, $r_2=g(I_2-\beta r_1)$. Додавши і віднявши ці рівності, отримаємо суму $r_1+r_2=g(I_1+I_2)/(1+g\beta)$ і різницю $r_1-r_2=g(I_1-I_2)/(1-g\beta)$. Їхнє відношення --- наскільки різко мережа віддає перевагу одному входу над іншим:
\begin{empheq}[box=\keybox]{equation}
\frac{r_1-r_2}{r_1+r_2} \;=\; \varepsilon\,\frac{1+g\beta}{1-g\beta}, \qquad \varepsilon=\frac{I_1-I_2}{I_1+I_2}, \quad g\beta<1 .
\label{eq:wtagain}
\end{empheq}
Мала різниця входів $\varepsilon$ підсилюється в $(1+g\beta)/(1-g\beta)$ раза, і цей множник необмежено зростає, коли $g\beta$ наближається до одиниці. При $g\beta>1$ стан, де працюють обидві групи, нестійкий, як у твердженні~\ref{prop:wta}: перемагає одна, інша мовчить (рис.~\ref{fig:pitchfork}).

\begin{figure}[t]
\centering
\begin{tikzpicture}[>=Stealth,x=3.2cm,y=1.5cm]
\draw[->,gray!70!black] (0,0) -- (2.15,0) node[below left,black] {\small $g\beta$};
\draw[->,gray!70!black] (0,-1.25) -- (0,1.35) node[above,black] {\small $(r_1-r_2)/(r_1+r_2)$};
\draw[gray!60!black] (1,-0.04) -- (1,0.04); \node[below right,font=\small] at (1,0) {$1$};
\node[left,font=\scriptsize] at (0,1) {$1$}; \node[left,font=\scriptsize] at (0,-1) {$-1$};
\draw[very thick,blue!60!black] plot[domain=0:0.905,samples=60] (\x,{0.05*(1+\x)/(1-\x)}) -- (1,1) -- (2.05,1);
\draw[very thick,blue!60!black] (1.105,-1) -- (2.05,-1);
\draw[thick,densely dashed,gray!60!black] (1,0) -- (2.05,0);
\node[font=\small,align=center] at (0.45,-0.62) {розмірковує:\\обидві групи\\працюють};
\node[font=\small,align=center] at (1.55,0.45) {вирішила:\\працює одна};
\node[font=\small,blue!60!black,anchor=south west] at (1.05,-0.97) {слабкий вхід переміг раніше --- тримається};
\end{tikzpicture}
\caption{Вибір із двох за різного коефіцієнта підсилення, входи відрізняються на $\varepsilon=0.05$. При $g\beta>1$ стійкі обидва рішення (суцільні лінії; перемога слабкого входу --- при $g\beta>(1+\varepsilon)/(1-\varepsilon)\approx1.1$), і мережа тримає те, яке вже ухвалила; симетричний стан (штрихова лінія) нестійкий.}
\label{fig:pitchfork}
\end{figure}

\begin{keyblock}
\begin{proposition}[Підсилення перемикає режим вибору]
\label{prop:gainwta}
У мережі вибору із взаємним гальмуванням $\beta$ коефіцієнт підсилення $g$ переводить мережу через межу $g\beta=1$. Нижче неї мережа «розмірковує»: обидві групи працюють, а різниця входів підсилюється за~\eqref{eq:wtagain}. Вище --- мережа «вирішила»: працює одна група, і рішення тримається. Змінюючи $g$, \nt{NORA} може перемикати мережу між цими режимами, не чіпаючи жодної ваги.
\end{proposition}
\end{keyblock}

\section{Підсилення як температура}

Тепер додамо до вибору шум, що виникає в самій мережі, після підсилювача. Яка група переможе, вирішує знак величини $g(u_A-u_B)+\eta$, де $u_A-u_B$ --- різниця входів, тобто різниця вивчених ваг із розділу~\ref{sec:dofa}, а $\eta$ --- шум із масштабом $\sigma$. Якщо шум розподілений за логістичним законом (для гаусового крива майже та сама), імовірність вибрати $A$ дорівнює
\begin{empheq}[box=\keybox]{equation}
P(A) \;=\; \frac{1}{1+e^{-g\,(u_A-u_B)/\sigma}} .
\label{eq:softmax}
\end{empheq}
Програміст упізнає тут вибір за softmax із температурою $T=\sigma/g$. Коефіцієнт підсилення --- обернена температура. За малого $g$ навіть помітно кращий варіант вибирається ненабагато частіше за гірший: мережа багато шукає. За великого $g$ найкращий відомий варіант вибирається майже завжди: мережа користується тим, що знає.

Уточнімо, що таке $g$ у~\eqref{eq:wtagain} і~\eqref{eq:softmax}: це \emph{діюче} підсилення різниці входів поточної задачі в момент вибору. Два режими \nt{NORA} діють на нього протилежно. Уривчаста пачка в момент рішення підвищує його, і мережа закріплює вибір. Висока фонова активність, навпаки, іде разом із пошуком: у людей перед вибором навмання базова зіниця ширша~\cite{JepmaNieuwenhuis2011}, а в щурів вхід \nt{NORA} у передню поясну кору переводить вибір у випадковий режим~\cite{Tervo2014stochastic}. Отже, високій фоновій \nt{NORA} відповідає \emph{мале} діюче $g$, а пачці --- велике.

\section{Скільки шукати}

Яке підсилення краще? Ми перевірили це на моделі. Мережа вибирає одну з двох дій за~\eqref{eq:softmax}; кожна приносить винагороду з певною ймовірністю, і мережа оцінює її за винагородами вибраної дії, як у розділі~\ref{sec:dofa}. Час від часу світ змінюється: обидві ймовірності розігруються заново. Змінитися може і вибрана дія, і та, яку мережа давно не пробувала; про другу мережа може дізнатися лише пошуком. На рис.~\ref{fig:invertedU} показано, яку частку найкращої можливої винагороди отримує мережа за різних сталих $g$.

\begin{figure}[t]
\centering
\begin{tikzpicture}[>=Stealth,x=1.2cm,y=1cm]
\draw[->,gray!70!black] (0,0) -- (7.6,0) node[below left,black] {\small $g$};
\draw[->,gray!70!black] (0,0) -- (0,4.4) node[above,black,font=\small] {частка найкращої винагороди};
\foreach \x/\l in {0/0.5,1/1,2/2,3/4,4/8,5/16,6/32,7/64}{\draw[gray!60!black] (\x,-0.06) -- (\x,0.06); \node[below,font=\scriptsize] at (\x,0) {\l};}
\foreach \v/\y in {0.8/0.8,0.9/2.4,1.0/4.0}{\draw[gray!60!black] (-0.06,\y) -- (0.06,\y); \node[left,font=\scriptsize] at (0,\y) {\v};}
\draw[very thick,blue!60!black] plot[mark=*,mark size=1.3pt] coordinates {(0,0.368) (1,0.880) (2,1.728) (3,2.784) (4,3.456) (5,3.616) (6,3.488) (7,3.264)};
\draw[very thick,red!70!black] plot[mark=*,mark size=1.3pt] coordinates {(0,0.368) (1,0.672) (2,1.184) (3,1.840) (4,2.176) (5,1.920) (6,1.456) (7,1.104)};
\node[blue!60!black,font=\small,anchor=south] at (5,3.7) {спокійний світ};
\node[red!70!black,font=\small,anchor=north] at (5.1,1.25) {світ, що швидко змінюється};
\draw[->,thick,blue!60!black] (5,3.25) -- (5,3.55);
\draw[->,thick,red!70!black] (4,1.8) -- (4,2.1);
\node[font=\small\itshape,gray!35!black,anchor=west] at (0.15,2.3) {шукає наосліп};
\node[font=\small\itshape,gray!35!black] at (6.45,2.55) {не помічає змін};
\end{tikzpicture}
\caption{Частка найкращої можливої винагороди за сталого коефіцієнта підсилення $g$ (шкала за $g$ логарифмічна). Синя крива --- світ змінюється в середньому раз на тисячу спроб, червона --- раз на двадцять. Стрілки позначають найкраще $g$: у світі, що швидко змінюється, воно вдвічі менше. Середнє за 60 прогонами по 4000 спроб.}
\label{fig:invertedU}
\end{figure}

При $g=0.5$ мережа майже не користується знанням і отримує близько трьох чвертей можливого, за дуже великого $g$ пропускає зміни. Найкраще значення: $g=16$, коли світ змінюється раз на тисячу спроб (частка $0.976$), і $g=8$, коли раз на двадцять ($0.886$).

\begin{keyblock}
\begin{proposition}[Перевернуте U]
\label{prop:explore}
Винагорода, яку збирає мережа, залежить від коефіцієнта підсилення як перевернуте U: надто мале $g$ витрачає спроби на пошук, надто велике не помічає змін. Найкраще $g$ тим менше, чим швидше змінюється світ.
\end{proposition}
\end{keyblock}

Перевернуте U спостерігається і в дослідах: тварини розв'язують задачу найкраще за помірної фонової активності \nt{NORA}-нейронів із чіткими уривчастими відповідями, а за високої фонової активності відволікаються і перемикаються між задачами~\cite{AstonJones2005}. Це узгоджується з різницею режимів із попереднього підрозділу: уривчаста пачка в момент рішення дає велике діюче $g$ саме тоді, коли треба вибрати, а висока фонова активність без пачок підсилює все підряд, зокрема й сторонні входи, тож діюче $g$ для поточної задачі падає, --- це пошук.

\section{Хто крутить ручку}

Якщо найкраще підсилення залежить від того, як швидко змінюється світ, його добре було б підлаштовувати. Але звідки \nt{NORA}-нейронам знати, що світ змінився? Природна ознака --- \emph{несподіваність}: помилки передбачення стали більшими, ніж зазвичай. Важливо розрізняти два види невизначеності. Якщо винагорода просто випадкова, помилки великі завжди, і це очікувано. Якщо ж помилки раптом зросли порівняно з їхньою звичайною величиною, отже, щось змінилося. За однією з гіпотез, саме цю неочікувану невизначеність і повідомляє \nt{NORA}, а очікувану --- \nt{ACET}~\cite{YuDayan2005}.

Що робити з цим сигналом? Можна знизити підсилення і більше шукати. Можна підвищити швидкість навчання, щоб швидше забути застарілі оцінки: досліди показують, що після різкої зміни люди навчаються швидше, і зіниця при цьому розширюється~\cite{Nassar2012}; нагадаємо, що зіниця --- показник не лише \nt{NORA}, а й \nt{ACET}~\cite{Reimer2016}. А можна зробити і те, й інше одразу: за ще однією гіпотезою, уривчаста пачка \nt{NORA} працює як \emph{переривання} --- сигнал, за яким мережа скидає поточний стан і перебудовується під нову обстановку~\cite{BouretSara2005}, як спільна лінія переривання процесора.

Чесно скажемо, чого ми не знайшли. У моделі ми пробували обидва прості способи підлаштування: знижувати $g$ або підвищувати швидкість навчання, коли згладжена помилка перевищує своє довготривале середнє. У світі, який то довго стоїть, то часто змінюється, найкращі налаштування обох способів дали $0.926$--$0.927$ від найкращої винагороди --- майже стільки ж, скільки найкраще стале $g$ ($0.925$ при $g=8$) або стала, але досить велика швидкість навчання ($0.928$), а невдалі налаштування --- менше. Криві на рис.~\ref{fig:invertedU} біля вершини пологі, і в такому світі підлаштовувати майже нічого. Підлаштування має окупатися там, де різні режими вимагають дуже різних налаштувань. Який саме закон керування реалізує \nt{NORA}, поки не встановлено.

\section{Підсумок}

\begin{table}[t]
\centering
\footnotesize
\setlength{\tabcolsep}{4pt}
\renewcommand{\arraystretch}{1.15}
\begin{tabular}{>{\raggedright}p{3.0cm}>{\raggedright}p{4.4cm}>{\raggedright\arraybackslash}p{6.0cm}}
\toprule
Що робить \nt{NORA} & Як & Що відомо \\
\midrule
змінює крутість відповіді нейронів & коефіцієнт підсилення $g$ у~\eqref{eq:gain} & підвищення відношення сигналу до шуму виміряно; мультиплікативна модель --- спрощення \\
перемикає вибір & переводить мережу через $g\beta=1$ & випливає з моделі~\eqref{eq:wtagain}; прямих вимірювань порога немає \\
задає, скільки шукати & $g$ --- обернена температура~\eqref{eq:softmax}; фон --- мале діюче $g$ (пошук), пачка --- велике (рішення) & перевернуте U і два режими роботи виміряно; зв'язок із пошуком --- добре обґрунтована гіпотеза \\
повідомляє про зміни & неочікувана невизначеність & зіниця і швидкість навчання зростають після змін --- виміряно; що це сигнал \nt{NORA} --- гіпотеза \\
перериває і скидає & уривчаста пачка на несподівану подію & пачки на важливі події виміряно; «переривання» --- гіпотеза \\
\bottomrule
\end{tabular}
\caption{Що додає \nt{NORA} до мережі з \nt{GLUT}, \nt{GABA} і \nt{DOPA}.}
\label{tab:nora}
\end{table}

\nt{DOPA} каже мережі, \emph{куди} змінювати ваги. \nt{NORA} не змінює жодної ваги, але вирішує, \emph{як} мережа користується тим, що вже знає (таблиця~\ref{tab:nora}): одна ручка --- коефіцієнт підсилення --- перетворює нейрон із підсилювача на компаратор, мережу вибору --- з тієї, що «розмірковує», на ту, що «вирішила», а вибір --- з пошуку на використання. Як \nt{NORA} дізнається, наскільки швидко змінюється світ, --- відкрите питання.

\section{Задачі}

\begin{exercise}[\lvlA]
\label{ex:08:1}
\emph{Одна ручка на весь мозок.} (а)~За таблицею~\ref{tab:types} оцініть, скільки нейронів мозку припадає на один \nt{NORA}-нейрон. (б)~Шум до підсилювача $\sigma_{\text{до}}=0.5$, після --- $\sigma_{\text{після}}=4$. За якого $g$ величина $d'$ з~\eqref{eq:gainsnr} досягає $90\,\%$ своєї границі?
\end{exercise}

\begin{exercise}[\lvlB]
\label{ex:08:2}
\emph{Вибір із коефіцієнтом підсилення.} $\tau\,\dd r_1/\dd t=-r_1+g[I_1-\beta r_2]_+$, $\tau\,\dd r_2/\dd t=-r_2+g[I_2-\beta r_1]_+$, $\tau=20\ms$. (а)~За який час загасає $r_1-r_2$ при $g\beta=0.5$ і $0.9$? (б)~При $\varepsilon=0.05$ знайдіть $g\beta$, нижче якого працюють обидві групи і вище якого стійка перемога слабкого входу.
\end{exercise}

\begin{exercise}[\lvlA]
\label{ex:08:3}
\emph{Softmax із шуму.} Кожна з $K$ груп отримує свій шум Гумбеля, $P(\eta_k<z)=\exp(-e^{-z/\sigma})$, і перемагає найбільше $gu_k+\eta_k$. Знайдіть $P(k)$. За якого $g$ кращий із двох варіантів з $u_A-u_B=0.1$, $\sigma=1$, вибирається в $90\,\%$ випадків?
\end{exercise}

\begin{exercise}[\lvlB]
\label{ex:08:4}
\emph{Оцінка найкращого підсилення.} У моделі рис.~\ref{fig:invertedU} імовірності винагород після зміни рівномірні на $[0,1]$. Мережа пробує гіршу дію раз на $e^{g\Delta}$ спроб; прирівнявши це до $1/h$, отримаємо $g^*\approx\ln(1/h)/\bar\Delta$. Знайдіть $\bar\Delta=\langle|p_A-p_B|\rangle$ і $g^*$ при $h=0.001$, $0.01$, $0.05$.
\end{exercise}

\begin{exercise}[\lvlB]
\label{ex:08:5}
\emph{Моделювання: підсилення і швидкість змін.} Знайдіть $g^*(h)$ при $h$ від $0.001$ до $0.2$ копією \texttt{models/\allowbreak nora\_explore.py} (вона пише файли в поточну теку). Де правильна оцінка задачі~\ref{ex:08:4} і чому за швидких змін $g^*$ перестає спадати?
\end{exercise}

\begin{exercise}[\lvlB]
\label{ex:08:6}
\emph{Проєктування: переривання для мережі вибору.} Мережа задачі~\ref{ex:08:2} з $\beta=2$, $\tau=20\ms$, $g=1$ тримає $r_1=0.9$, $r_2=0$, хоча входи тепер $I_1=0.9$, $I_2=1.0$. \nt{NORA} на час $T_p$ знижує $g$ до $g_{\text{lo}}$. Знайдіть найменше $T_p$ при $g_{\text{lo}}=0$ аналітично і моделлю при $g_{\text{lo}}=0.25$.
\end{exercise}

\begin{exercise}[\lvlC]
\label{ex:08:7}
\emph{Коли підлаштування окупається.} Оракул знає, спокійна епоха чи мінлива, і ставить у кожній своє $g$. (а)~Моделлю знайдіть його виграш над найкращим сталим $g$ у світі розділу (епохи по $1000$ спроб, $h=0.001$ і $0.05$). (б)~Побудуйте світ, де виграш більший, і знайдіть, яку його частину забирає підлаштування за несподіваністю.
\end{exercise}
